\chapter{Tudo o Que Surge Acaba}

\lettrine{O}{\space Buda disse que} a origem de todo sofrimento é a ignorância -- por isso é
importante considerar o que ele realmente quis dizer por ``ignorância''. Na sua
maioria, os seres humanos no mundo vivem como se realmente fossem os seus
próprios hábitos, pensamentos, sentimentos e memórias. Não usam o tempo ou não
têm a oportunidade de olhar para as suas vidas, de observarem e considerarem
como estas condições funcionam.

O que é uma condição? O corpo com o qual estamos, as emoções e sentimentos, as
percepções da mente, concepções e consciência através dos sentidos -- estas são
as condições. Uma condição é algo que é adicionado e composto; algo que surge e
acaba; não é o incriado, não nascido, a realidade última não originada.

A religião é o que os seres humanos usam para tentarem regressar à realização
última para além dos ciclos de nascimento e morte, a sabedoria supramundana ou
\emph{lokuttara paññā}; Nirvāna ou Nibbāna é a experiência dessa
realidade transcendente. Isto é quando subitamente conhecemos a verdade, não por
estudar as escrituras Pāli ou um livro Zen, mas por experiência directa.

Geralmente concebemos a verdade como sendo \emph{algo} e Nibbāna como
sendo um estado mental pacífico ou experiência extática. Todos nós já
experimentámos algum tipo de felicidade, então gostamos de conceber o Não
Nascido, Incriado, Não Originado como uma experiência feliz. Mas o Buddha foi
muito cuidadoso em nunca descrever a Realidade Suprema ou Nibbāna -- ele
nunca falou muito sobre isso. As pessoas querem saber o que é, escrever livros
sobre isso e especular sobre a natureza do Nibbāna -- mas isso é
exactamente o que o Buddha não fez.

Em vez disso, ele apontou para o conhecimento directo das condições que mudam,
aquilo que podemos saber pela nossa própria experiência neste momento. Não se
trata de acreditar em qualquer outro alguém. É uma questão de saber, neste
momento presente, que tudo o que surge acaba.

Portanto, colocamos esse tipo de atenção nas nossas vidas -- para estarmos
atentos e percebermos que tudo o que surge acaba; que qualquer condição da nossa
mente ou corpo -- seja uma sensação de prazer ou dor, sentimento ou memória,
visão, som, cheiro, paladar ou tacto, dentro ou fora -- é apenas uma condição.

É importante reflectir sobre o significado real de ``ignorância'' no sentido em
que o Buddha disse quando lhe chamou a origem de todo sofrimento. ``Ser
ignorante'' significa que nos identificamos com essas condições, considerando-as
como ``eu'' ou ``meu'' ou como algo que não queremos ser ``eu'' ou ``meu''.
Temos a ideia de que temos que encontrar alguma condição agradável permanente,
conseguir algo, obter algo que não temos. Mas podemos perceber que o desejo na
mente é uma coisa em movimento; está à procura de algo, portanto é uma condição
em mudança que surge e acaba -- é não-eu (não é eu). A expressão ``não-eu''
(\emph{anatta}) não é um tipo qualquer de mantra\footnote{Uma palavra ou frase
  com significado espiritual. É usado como um foco de atenção e reflexão ao se
  repetir muitas vezes na meditação.} que usamos para nos livrarmos das coisas,
mas é uma penetração real compreensiva da natureza básica de todos os desejos.

Ao olhar com cuidado, com muita paciência e humildade, começamos a ver que o
criado surge do Incriado e volta ao Incriado; desaparece e nada mais resta. Se
fosse realmente eu e realmente meu, ficaria, certo? Se fosse realmente meu, para
onde iria, para algum tipo de arrecadação de personalidade? Mas esse conceito e
tudo o que se possa conceber é uma condição que surge e acaba. Sempre que nos
tentamos conceber, qualquer conceito ou memória de nós como isto ou aquilo é
apenas uma condição da nossa mente. Não é o que somos -- nós não somos uma
condição da nossa mente. Então, tristeza, desespero, amor e felicidade são todas
condições da mente e são todas não-eu.

Observemos na nossa vida quando sofremos ou sentimos descontentamento -- porquê?
É por causa de algum apego, alguma ideia de nós ou de outra pessoa. Alguém que
amamos morre e sentimos pena de nós mesmos; pensamos nos bons momentos que
tivemos e concentramo-nos nisso, criando mais condições da mente. Talvez
sintamos culpa por não termos dado ou amado o tempo inteiro -- sendo esta também
uma condição da mente. Temos uma memória; concebemo-los como estando vivos --
mas a própria ideia de uma pessoa é uma percepção da mente -- não uma pessoa, ou
é? Lembramo-nos de alguém ainda vivo, com quem gostaríamos de estar agora -- essa
é uma condição da mente; ou lembramo-nos de alguém que morreu, que jamais
veremos -- isso também é uma condição da mente.

A meditação budista é uma forma de ver as condições da mente, investigando e
vendo o que são, em vez de acreditar nelas. As pessoas querem acreditar --
quando alguém próximo de nós morre, alguém acaba por nos dizer: ``Oh, subiu ao
céu com Deus Pai'', ou ``estão agora a viver nas delícias do céu Tusita.'' Dizem
isto de forma a termos uma percepção agradável mental -- ``Bem, agora sei que a
minha avó está feliz lá em cima nos reinos celestiais, dançando com os anjos''.
Então, outro alguém diz ``Bem, sabe, ela fez algumas coisas bem más, ela
provavelmente está lá em baixo no Inferno, ardendo nos fogos eternos!'' Com
isto, logo nos começamos a preocupar que igualmente lá acabemos talvez -- mas
isso é uma percepção da mente. Céu e Inferno são fenómenos condicionados. Então
-- se reflectirmos sobre algo sucedido há dez anos\ldots{} essa é uma condição da
mente que surge e acaba, e a razão porque surge é porque eu acabei de a sugerir.
Portanto, essa condição depende de outra condição, a memória é o que
experienciámos e o futuro é o desconhecido.

Mas quem é que conhece as condições do momento? Não o consigo encontrar: só o
saber existe, e o saber pode conhecer tudo o que está presente agora --
agradável ou desagradável -- especulações sobre o futuro ou reminiscências do
passado -- criações de nós mesmos como isto ou aquilo. Criamo-nos a nós mesmos
ou ao mundo em que vivemos -- então não podemos realmente culpar mais ninguém.
Se o fazemos, é porque ainda somos ignorantes. Denominamos ``Aquele que Sabe''
de ``Buddha'' -- mas isso não significa que ``Buddha'' seja uma condição. Não é
dizer que esta Buddha-rupa (estátua de Buddha) saiba alguma coisa; mais do que
isso, significa que ``Buddha'' é o saber. Portanto, a meditação budista é
realmente estar ciente, mais do que se tornar Buddha.

A ideia de nos tornarmos Buddha é baseada em condições -- pensamos ser alguém
que agora não é Buddha, e de forma a nos tornarmos Buddha temos que ler livros
para descobrir como nos tornaremos num. Claro, isso significa que temos que
trabalhar muito para nos livrarmos das qualidades que não são características de
Buddha; estamos longe de ser perfeitos, irritamo-nos, somos gananciosos,
duvidosos e amedrontados, e claro, os Buddhas não têm isto -- porque Buddha é
aquilo que sabe, portanto eles sabem melhor. Então, para nos tornarmos Buddha,
temos que nos livrar dessas coisas que não são características de Buddha e
tentar obter qualidades características de Buddha como compaixão e todos estes
tipos de coisas. \emph{E tudo isto são criações da mente!} Portanto, criamos
``Buddhas'' porque acreditamos nas criações da mente. Mas elas não são Buddhas
reais. Elas são apenas falsos Buddhas. Eles não são Buddhas de sabedoria, são
apenas condições da nossa mente.

Enquanto nos concebermos como alguém que tem que fazer algo para se tornar outra
coisa, ainda caímos numa armadilha, uma condição da mente como sendo um eu, e
nunca entenderemos nada como deve ser. Independentemente de quantos anos
meditemos, nunca realmente entenderemos o ensinamento; simplesmente falharemos
sempre o alvo. A maneira directa de ver as coisas \emph{agora} -- de que tudo o
que surge acaba -- não significa que estamos a desprezar o que seja. Significa
que estamos a observar de uma maneira que nunca nos preocupámos em observar
antes. Estamos a observar de uma perspectiva daquilo que está aqui e agora, em
vez de procurar por algo que não está aqui. Então, se entramos na Sala de
Meditação a pensar: ``Eu tenho que passar esta hora a procurar pelo Buddha,
tentar tornar-me algo, tentar livrar-me destes pensamentos nefastos, sentar e
praticar duro, tentar tornar-me o que eu me deveria tornar -- então vou
sentar-me aqui e tentar livrar-me de coisas, tentar obter coisas, tentar segurar
as coisas'' \ldots{} com essa atitude, a meditação é um esforço realmente
extenuante e sempre um fracasso.

Mas, se em vez disso, entramos na Sala de Meditação e estamos simplesmente
cientes das condições da mente, vemos em perspectiva o desejo de se tornar, de
se livrar de, de fazer algo ou a sensação de que não o conseguimos fazer; ou que
se é um especialista, seja o que for -- começamos a ver que tudo o que estamos a
experienciar é uma condição que muda e não ``eu''. Observamos a perspectiva de
ser Buddha, em vez de fazer algo para nos tornarmos Buddha. Quando falamos sobre
\emph{sati}, consciência, é isto que queremos dizer.

Estou realmente chocado e surpreso com muitas pessoas religiosas -- cristãos ou
budistas ou o que seja -- que parecem ignorar a prática da sua religião. Poucas
pessoas parecem ter qualquer perspectiva sobre a doutrina religiosa, crença e
descrença. Não se preocupam em descobrir. Ainda estão a tentar descrever o
indescritível, limitar o ilimitado, conhecer o incognoscível\ldots{} e poucas
pessoas reparam na forma como se encontram. Acreditam naquilo que outro alguém
lhes disse.

Hoje em dia, monges Budistas Theravāda dir-vos-ão que vocês não se podem
iluminar, já nem podemos sequer alcançar a entrada-no-caminho (corrente), o
primeiro estágio da santidade, que esses tempos são passado. Eles acreditam que
a iluminação é uma possibilidade tão remota que já nem eles próprios sequer
colocam muito esforço para verem que tudo o que surge acaba. Portanto, os monges
podem passar vidas inteiras a ler livros e a traduzir Suttas, acreditando ainda
que não são iluminados e que é impossível sê-lo. Mas, então, qual o sentido da
religião? Porquê então preocupar-nos se a verdade última é assim tão remota e
uma possibilidade tão improvável? Tornamo-nos simplesmente como antropólogos,
sociólogos ou filósofos discutindo religião comparada.

Gotama, o Buddha, era alguém cuja sabedoria resultou de observar a Natureza, de
observar as condições da mente e do corpo. Ora, isso não é impossível de fazer
para qualquer um de nós. Temos mentes e corpos; tudo o que precisamos de fazer é
observá-los. Não é como se tivéssemos que ter poderes especiais para fazer isso
ou, que este tempo actual seja de alguma forma diferente do de Gotama, o Buddha.
O tempo é uma ilusão causada pela ignorância. As pessoas da época de Gotama, o
Buddha, não eram assim tão diferentes de agora -- elas tinham ganância, ódio e
ilusão, egos, presunções e medos exactamente como as pessoas de hoje. Se
começarmos a pensar sobre as doutrinas budistas e os diferentes níveis de
realização, entraremos simplesmente em estado de dúvida. Não precisamos de
verificar em nós próprios como numa lista de um livro -- mas sim sabermos por
nós mesmos até que nenhuma condição do corpo ou da mente nos iluda.

As pessoas dizem-me: ``Eu não consigo fazer tudo isso. Sou uma pessoa comum, um
leigo apenas; quando penso em fazer tudo isso, percebo que é demasiado para mim,
não consigo fazer isso.'' Eu digo, ``Se pensar nisso, não o conseguirá fazer, é
simples. Não pense nisso, faça-o simplesmente.'' O pensamento apenas o leva à
dúvida. As pessoas que pensam acerca da vida não conseguem fazer nada. Se vale a
pena fazer, faça. Quando ficar deprimido, aprenda com a depressão; quando ficar
doente, aprenda com a doença; quando estiver feliz, aprenda com a felicidade --
todas estas são oportunidades de aprender no mundo. Continue a ouvir e a
observar em silêncio como um modo de vida\ldots{} então começará a compreender
as condições. Não há nada a temer. Não há nada que precise de obter que não
tenha, não existe nada para abandonar.

% FIXME ordenação => taking precepts

\photoPage{image-003-anagarika-line.png}{Anagārikas na sua cerimónia de ordenação}{}
